% ---- score2 standard preamble (from: python3 score2_tools.py preamble) ----
\documentclass[10pt]{extarticle}
\usepackage[a4paper,margin=1.5cm]{geometry}
\usepackage[T1]{fontenc}
\usepackage{amsmath,amssymb}
\usepackage{graphicx}
\usepackage[export]{adjustbox}
\usepackage{multicol}
\usepackage{enumitem}
\usepackage{booktabs,array,tabularx}
\setlength{\columnsep}{0.7cm}
\setlength{\parindent}{0pt}
\setlength{\parskip}{3pt}
\pagestyle{empty}
\makeatletter
% figure registry (figdefs.tex is generated by the extract step)
\newcommand{\qfigdef}[4]{\@namedef{qf@file@#1}{#2}\@namedef{qf@w@#1}{#3}\@namedef{qf@wide@#1}{#4}}
\newcommand{\qf@check}[1]{\@ifundefined{qf@file@#1}{\PackageError{score2}{Unknown figure id #1}{}}{}}
\newcommand{\qf@img}[4][]{\includegraphics[width=#2\dimexpr\@nameuse{qf@w@#4}pt\relax,max width=#3,max height=0.45\textheight,#1]{\@nameuse{qf@file@#4}}}
% Full-width questions: finalize wraps every question that contains a wide figure in \qfullbegin...\qfullend,
% so the WHOLE question (stem, figures, options) is set across the page in reading order, then the two columns
% resume. The question's head - its number, stem and first figure (\qfullsplit marks the end of the head) -
% is kept on one page, so a number is never left at a page foot with its figure overleaf; the rest may break
% across pages normally, so no big gaps are left. Without a \qfullsplit, the whole question is kept together
% if it is under 45% of a page. Never write these yourself.
\newif\ifqf@full
\newif\ifqf@boxed
\newsavebox{\qf@box}
\newcommand{\qf@open}{\global\qf@fulltrue\par\global\qf@boxedtrue
  \setbox\qf@box\vbox\bgroup\hsize\textwidth\linewidth\textwidth}
\newcommand{\qf@emit}[1]{\par\egroup\global\qf@boxedfalse
  \ifdim\dimexpr\ht\qf@box+\dp\qf@box\relax<#1\textheight\noindent\usebox{\qf@box}\par\else\unvbox\qf@box\fi}
\newcommand{\qfullsplit}{\ifqf@boxed\qf@emit{0.85}\fi}
\newcommand{\qf@close}{\ifqf@boxed\qf@emit{0.45}\fi\global\qf@fullfalse}
\newcommand{\qfullbegin}{\end{multicols}\qf@open}
\newcommand{\qfullend}{\qf@close\begin{multicols}{2}}
\newcommand{\qfullstart}{\qf@open}
\newcommand{\qfulldone}{\qf@close}
% \qfig[scale]{id}: block figure, centred, directly where it is written (size set by the tool).
% Fallback if a wide figure is not inside a full-width question: break out of the columns in place - no floats.
\newcommand{\qfig}[2][1]{\qf@check{#2}%
  \ifqf@full\qf@blk{#1}{#2}\else
  \ifnum\@nameuse{qf@wide@#2}=1 \expandafter\@firstoftwo\else\expandafter\@secondoftwo\fi
  {\end{multicols}\begin{center}\qf@img{#1}{\textwidth}{#2}\end{center}\begin{multicols}{2}}%
  {\qf@blk{#1}{#2}}\fi}
\newcommand{\qf@blk}[2]{\par\smallskip{\centering\qf@img{#1}{\linewidth}{#2}\par}\smallskip}
% \qfigi[scale]{id}: inline figure (inside option lists / tables / minipages); never breaks columns.
\newcommand{\qfigi}[2][1]{\qf@check{#2}\qf@img[valign=c]{#1}{\linewidth}{#2}}
\makeatother
% structure macros - lint reads these to check completeness and figure placement
\newcommand{\papersec}[1]{\par\bigskip{\centering\Large\bfseries PAPER-#1\par}\medskip}
\newcommand{\partsec}[2]{\par\medskip\noindent{\large\bfseries PART-#1: #2}\par\smallskip}
\newcommand{\sect}[1]{\par\smallskip\noindent\textbf{#1}\par}
\newcommand{\qn}[1]{\par\smallskip\noindent\textbf{#1.}~}
\newcommand{\paraq}[2]{\par\smallskip\noindent\textbf{Paragraph for Questions #1 and #2}\par}
\newenvironment{qoptions}{\begin{itemize}[leftmargin=1.8em,itemsep=1pt,topsep=2pt,label={}]}{\end{itemize}}
\newcommand{\opt}[2]{\item[(#1)] #2}
% figure registry (generated by score2_tools.py finalize)
\qfigdef{p1_phy_s2_q05_1}{Pre-Score (2023-24) - Test 06 (20 Nov 2023)-img/p1_phy_s2_q05_1.png}{134.7}{0}
\qfigdef{p1_phy_s2_q07_1}{Pre-Score (2023-24) - Test 06 (20 Nov 2023)-img/p1_phy_s2_q07_1.png}{74.2}{0}
\qfigdef{p1_phy_s4_q04_1}{Pre-Score (2023-24) - Test 06 (20 Nov 2023)-img/p1_phy_s4_q04_1.png}{92.6}{0}
\qfigdef{p1_phy_s4_q05_1}{Pre-Score (2023-24) - Test 06 (20 Nov 2023)-img/p1_phy_s4_q05_1.png}{144.2}{0}
\qfigdef{p1_phy_s4_q06_1}{Pre-Score (2023-24) - Test 06 (20 Nov 2023)-img/p1_phy_s4_q06_1.png}{192.8}{0}
\qfigdef{p1_chem_s1_q01_1}{Pre-Score (2023-24) - Test 06 (20 Nov 2023)-img/p1_chem_s1_q01_1.png}{118.5}{0}
\qfigdef{p1_chem_s2_q05_A}{Pre-Score (2023-24) - Test 06 (20 Nov 2023)-img/p1_chem_s2_q05_A.png}{65.6}{0}
\qfigdef{p1_chem_s2_q05_B}{Pre-Score (2023-24) - Test 06 (20 Nov 2023)-img/p1_chem_s2_q05_B.png}{67.1}{0}
\qfigdef{p1_chem_s2_q05_C}{Pre-Score (2023-24) - Test 06 (20 Nov 2023)-img/p1_chem_s2_q05_C.png}{62.9}{0}
\qfigdef{p1_chem_s2_q05_D}{Pre-Score (2023-24) - Test 06 (20 Nov 2023)-img/p1_chem_s2_q05_D.png}{47.7}{0}
\qfigdef{p1_chem_s3_q09_1}{Pre-Score (2023-24) - Test 06 (20 Nov 2023)-img/p1_chem_s3_q09_1.png}{52.0}{0}
\qfigdef{p1_chem_s3_q09_2}{Pre-Score (2023-24) - Test 06 (20 Nov 2023)-img/p1_chem_s3_q09_2.png}{33.7}{0}
\qfigdef{p1_chem_s3_q09_3}{Pre-Score (2023-24) - Test 06 (20 Nov 2023)-img/p1_chem_s3_q09_3.png}{41.9}{0}
\qfigdef{p1_chem_s3_q10_1}{Pre-Score (2023-24) - Test 06 (20 Nov 2023)-img/p1_chem_s3_q10_1.png}{100.6}{0}
\qfigdef{p1_chem_s3_q10_2}{Pre-Score (2023-24) - Test 06 (20 Nov 2023)-img/p1_chem_s3_q10_2.png}{78.2}{0}
\qfigdef{p1_chem_s3_q10_3}{Pre-Score (2023-24) - Test 06 (20 Nov 2023)-img/p1_chem_s3_q10_3.png}{62.5}{0}
\qfigdef{p1_chem_s3_q10_4}{Pre-Score (2023-24) - Test 06 (20 Nov 2023)-img/p1_chem_s3_q10_4.png}{107.0}{0}
\qfigdef{p1_chem_s4_q01_1}{Pre-Score (2023-24) - Test 06 (20 Nov 2023)-img/p1_chem_s4_q01_1.png}{271.8}{1}
\qfigdef{p1_chem_s4_q01_2}{Pre-Score (2023-24) - Test 06 (20 Nov 2023)-img/p1_chem_s4_q01_2.png}{245.2}{0}
\qfigdef{p1_chem_s4_q03_1}{Pre-Score (2023-24) - Test 06 (20 Nov 2023)-img/p1_chem_s4_q03_1.png}{276.8}{1}
\qfigdef{p1_chem_s4_q04_1}{Pre-Score (2023-24) - Test 06 (20 Nov 2023)-img/p1_chem_s4_q04_1.png}{245.2}{0}
\qfigdef{p1_chem_s4_q04_2}{Pre-Score (2023-24) - Test 06 (20 Nov 2023)-img/p1_chem_s4_q04_2.png}{239.5}{0}
\qfigdef{p1_chem_s4_q05_1}{Pre-Score (2023-24) - Test 06 (20 Nov 2023)-img/p1_chem_s4_q05_1.png}{170.7}{0}
\qfigdef{p1_chem_s4_q05_2}{Pre-Score (2023-24) - Test 06 (20 Nov 2023)-img/p1_chem_s4_q05_2.png}{150.9}{0}
\qfigdef{p1_chem_s4_q05_3}{Pre-Score (2023-24) - Test 06 (20 Nov 2023)-img/p1_chem_s4_q05_3.png}{166.0}{0}
\qfigdef{p1_chem_s4_q06_1}{Pre-Score (2023-24) - Test 06 (20 Nov 2023)-img/p1_chem_s4_q06_1.png}{212.0}{0}
\qfigdef{p2_phy_s1_q04_1}{Pre-Score (2023-24) - Test 06 (20 Nov 2023)-img/p2_phy_s1_q04_1.png}{116.3}{0}
\qfigdef{p2_phy_s2_q05_1}{Pre-Score (2023-24) - Test 06 (20 Nov 2023)-img/p2_phy_s2_q05_1.png}{113.7}{0}
\qfigdef{p2_phy_s2_q07_1}{Pre-Score (2023-24) - Test 06 (20 Nov 2023)-img/p2_phy_s2_q07_1.png}{108.3}{0}
\qfigdef{p2_phy_s3_para1_2_1}{Pre-Score (2023-24) - Test 06 (20 Nov 2023)-img/p2_phy_s3_para1_2_1.png}{74.2}{0}
\qfigdef{p2_phy_s3_para3_4_1}{Pre-Score (2023-24) - Test 06 (20 Nov 2023)-img/p2_phy_s3_para3_4_1.png}{105.6}{0}
\qfigdef{p2_phy_s4_q06_1}{Pre-Score (2023-24) - Test 06 (20 Nov 2023)-img/p2_phy_s4_q06_1.png}{150.1}{0}
\qfigdef{p2_phy_s4_q07_1}{Pre-Score (2023-24) - Test 06 (20 Nov 2023)-img/p2_phy_s4_q07_1.png}{89.9}{0}
\qfigdef{p2_chem_s1_q01_1}{Pre-Score (2023-24) - Test 06 (20 Nov 2023)-img/p2_chem_s1_q01_1.png}{112.0}{0}
\qfigdef{p2_chem_s1_q02_1}{Pre-Score (2023-24) - Test 06 (20 Nov 2023)-img/p2_chem_s1_q02_1.png}{206.5}{0}
\qfigdef{p2_chem_s1_q03_A}{Pre-Score (2023-24) - Test 06 (20 Nov 2023)-img/p2_chem_s1_q03_A.png}{89.1}{0}
\qfigdef{p2_chem_s1_q03_B}{Pre-Score (2023-24) - Test 06 (20 Nov 2023)-img/p2_chem_s1_q03_B.png}{104.8}{0}
\qfigdef{p2_chem_s1_q03_C}{Pre-Score (2023-24) - Test 06 (20 Nov 2023)-img/p2_chem_s1_q03_C.png}{129.2}{0}
\qfigdef{p2_chem_s1_q04_1}{Pre-Score (2023-24) - Test 06 (20 Nov 2023)-img/p2_chem_s1_q04_1.png}{54.8}{0}
\qfigdef{p2_chem_s2_q05_B}{Pre-Score (2023-24) - Test 06 (20 Nov 2023)-img/p2_chem_s2_q05_B.png}{76.8}{0}
\qfigdef{p2_chem_s2_q06_A}{Pre-Score (2023-24) - Test 06 (20 Nov 2023)-img/p2_chem_s2_q06_A.png}{131.0}{0}
\qfigdef{p2_chem_s2_q06_B}{Pre-Score (2023-24) - Test 06 (20 Nov 2023)-img/p2_chem_s2_q06_B.png}{111.9}{0}
\qfigdef{p2_chem_s2_q06_C}{Pre-Score (2023-24) - Test 06 (20 Nov 2023)-img/p2_chem_s2_q06_C.png}{74.5}{0}
\qfigdef{p2_chem_s2_q06_D}{Pre-Score (2023-24) - Test 06 (20 Nov 2023)-img/p2_chem_s2_q06_D.png}{99.0}{0}
\qfigdef{p2_chem_s2_q07_1}{Pre-Score (2023-24) - Test 06 (20 Nov 2023)-img/p2_chem_s2_q07_1.png}{292.4}{1}
\qfigdef{p2_chem_s3_q01_1}{Pre-Score (2023-24) - Test 06 (20 Nov 2023)-img/p2_chem_s3_q01_1.png}{31.6}{0}
\qfigdef{p2_chem_s3_q03_1}{Pre-Score (2023-24) - Test 06 (20 Nov 2023)-img/p2_chem_s3_q03_1.png}{108.3}{0}
\qfigdef{p2_chem_s4_q05_1}{Pre-Score (2023-24) - Test 06 (20 Nov 2023)-img/p2_chem_s4_q05_1.png}{264.5}{1}
\qfigdef{p2_chem_s4_q06_1}{Pre-Score (2023-24) - Test 06 (20 Nov 2023)-img/p2_chem_s4_q06_1.png}{124.4}{0}
% ---- body skeleton ----
% \begin{document}
% \begin{center}{\Large\bfseries JEE (Main+Advanced) Score-II --- Test NN (DD Mon YYYY)}\end{center}
% \begin{multicols}{2}
%   \papersec{1}   (only for two-paper tests)
%   \partsec{1}{PHYSICS} \sect{Section-I (i): Single Correct} \qn{1} ... \qfig{p1_phy_q01_1} ...
% \end{multicols}
% \end{document}

\begin{document}
\begin{center}{\Large\bfseries Pre-Score (2023-24) --- Test 06 (20 Nov 2023)}\end{center}\begin{multicols}{2}
\papersec{1}
\partsec{1}{PHYSICS}
\sect{SECTION-I(i)}
\qn{1} Ishan wants to float a ball (perfectly reflecting) in the air. He turned on a monochromatic light source at the bottom, perpendicular to the ground, pointing upwards with intensity $P$. The ball is dropped from a height $H$, directly into the light beam. The ball has a radius $R$ and a mass $m$. Unfortunately, Ishan did not compute the needed intensity correctly, and the light beam cannot keep it floating in the air. At what speed does the ball hit the ground? (Symbols have their usual meaning.)
\begin{qoptions}
\opt{A}{$\sqrt{H\left(g-\dfrac{\pi R^2P}{mc}\right)}$}
\opt{B}{$\sqrt{2H\left(g-\dfrac{\pi R^2P}{2mc}\right)}$}
\opt{C}{$2\sqrt{H\left(g-\dfrac{\pi R^2P}{mc}\right)}$}
\opt{D}{$\sqrt{2H\left(g-\dfrac{\pi R^2P}{mc}\right)}$}
\end{qoptions}
\qn{2} Gill has a dry well with constant circular cross-section and depth $h$. When standing at the bottom, he is able to see an angle $2\alpha\le 90^\circ$ of the sky above him. We want to increase this angle by pouring some specific liquid into the well (Gill still stay at the bottom). Find the condition for the index of refraction of the liquid if we want the angle of the visible sky to be twice as large.
\begin{qoptions}
\opt{A}{$n\ge 2\cos\alpha$}
\opt{B}{$n\le 2\cos\alpha$}
\opt{C}{$n\ge 2\sin\alpha$}
\opt{D}{$n\le 2\sin\alpha$}
\end{qoptions}
\qn{3} Rohit wants to find out how the refractive index of the atmosphere should depend on the height $y$ above the Earth's surface, if rays in the atmosphere were to travel along parabolic trajectories $y=x^2$. The dependence of refractive index on $y$ may be (You may assume earth to be flat):
\begin{qoptions}
\opt{A}{$n=n_0\sqrt{2y+1}$}
\opt{B}{$n=n_0\sqrt{4y+1}$}
\opt{C}{$n=n_0\sqrt{16y+1}$}
\opt{D}{None}
\end{qoptions}
\qn{4} The energy--momentum relation is the relativistic equation relating total energy (which is also called relativistic energy) to invariant mass (which is also called rest mass) and momentum. It is the extension of mass--energy equivalence for bodies or systems with nonzero momentum. It can be written as the following equation $E=\sqrt{(m_0c^2)^2+(pc)^2}=mc^2$. This equation holds for a body or system, such as one or more particles, with total energy $E$, invariant mass $m_0$, and momentum of magnitude $p$; the constant $c$ is the speed of light. X-ray photons are scattered through an angle of $90^\circ$ by electrons initially at rest. What is the change in the wavelength of the photons? (Take rest mass of electron as $m_0$)
\begin{qoptions}
\opt{A}{$\dfrac{h}{\sqrt2\,m_0c}$}
\opt{B}{$\dfrac{\sqrt2\,h}{m_0c}$}
\opt{C}{$\dfrac{h}{m_0c}$}
\opt{D}{None}
\end{qoptions}
\sect{SECTION-I(ii)}
\qn{5} You are given a long, thin polished metal tube and asked to put your eye at one end and to point the other end at a diffuse light source (a well-illuminated wall, for example). What you see is a pattern of rings. At the centre there is a bright circle, which is the diffuse light source seen directly through the hole in the end of the tube. Surrounding this is a ring of light that is less bright than the central circle. Around that there's another ring that's less bright still. And so on. The rings become successively less bright until the pipe appears dark. The rings appear to have approximately the same width as they get larger. The angle subtended at the eye by ring may be : (Diameter of tube is $D$, length of tube is $L$)
\qfig{p1_phy_s2_q05_1}
\begin{qoptions}
\opt{A}{$\dfrac{D}{2L}$}
\opt{B}{$\dfrac{D}{L}$}
\opt{C}{$\dfrac{3D}{2L}$}
\opt{D}{$\dfrac{2D}{L}$}
\end{qoptions}
\qn{6} What atomic number would an element have to have for the ground state speed of an electron to exceed the speed of light, if we do not consider relativistic corrections? (Use the Bohr model of the atom and assume that the atom is ionized to exactly one electron.)
\begin{qoptions}
\opt{A}{108}
\opt{B}{118}
\opt{C}{128}
\opt{D}{138}
\end{qoptions}
\qn{7} On a big room's wall there is a mirror of width 2 m. A person is standing 2 m away from the mirror and 1 m away from the wall. The person starts to move parallel to the mirror with a speed of 1 m/s. At the same moment the person's acquaintance starts to move with a speed of 1 m/s towards the mirror on the mirror's centerline. Initially, the acquaintance is standing at a distance 3.5 m from the mirror. After what time do the two see each other from the mirror?
\qfig{p1_phy_s2_q07_1}
\begin{qoptions}
\opt{A}{1 sec}
\opt{B}{1.5 sec}
\opt{C}{2 sec}
\opt{D}{2.5 sec}
\end{qoptions}
\sect{SECTION-I(iii)}
\qn{8} List-I contain hydrogen like atom/ion and List-II contains their properties according to Bohr's model. Select the correct match.
\begin{center}\footnotesize\renewcommand{\arraystretch}{1.6}\setlength{\tabcolsep}{3pt}
\begin{tabular}{|l@{\hspace{3pt}}p{2.5cm}|l@{\hspace{3pt}}p{3.2cm}|}
\hline \multicolumn{2}{|c|}{List-I} & \multicolumn{2}{c|}{List-II} \\\hline
(P) & H, electron in ground state. & (1) & Number of revolutions per seconds is $\dfrac{me^4}{\varepsilon_0^2h^3}$. \\\hline
(Q) & He$^+$, electron in ground state. & (2) & Radius of orbit is $\dfrac{\varepsilon_0h^2}{\pi me^2}$. \\\hline
(R) & Li$^{2+}$, electron in $3^{\text{rd}}$ orbit. & (3) & Velocity of electron is $\dfrac{e^2}{\varepsilon_0h}$. \\\hline
(S) & Be$^{3+}$, electron in $2^{\text{nd}}$ orbit. & (4) & Kinetic energy of electron is $\dfrac{me^4}{8\varepsilon_0^2h^2}$. \\\hline
\multicolumn{2}{|c|}{} & (5) & Potential energy is $-\dfrac{me^4}{\varepsilon_0^2h^2}$. \\\hline
\end{tabular}
\end{center}
\begin{qoptions}
\opt{A}{P $\to$ 4; Q $\to$ 2; R $\to$ 1; S $\to$ 3}
\opt{B}{P $\to$ 2; Q $\to$ 4; R $\to$ 1; S $\to$ 5}
\opt{C}{P $\to$ 2; Q $\to$ 3; R $\to$ 4; S $\to$ 5}
\opt{D}{P $\to$ 2; Q $\to$ 4; R $\to$ 2; S $\to$ 3}
\end{qoptions}
\qn{9} Select the correct match.
\begin{center}\footnotesize\renewcommand{\arraystretch}{1.3}\setlength{\tabcolsep}{3pt}
\begin{tabular}{|l@{\hspace{3pt}}p{2.6cm}|l@{\hspace{3pt}}p{2.6cm}|}
\hline \multicolumn{2}{|c|}{List-I} & \multicolumn{2}{c|}{List-II} \\\hline
(P) & The filament voltage in X-ray experiment is increased. & (1) & The average kinetic energy of the electrons decreases. \\\hline
(Q) & In photoelectric effect work function of the target is increased. & (2) & The average kinetic energy of the electron increases. \\\hline
(R) & The de-Broglie wavelength of $\beta$-rays is increased. & (3) & The intensity increases. \\\hline
(S) & The de-Broglie wavelength of $\beta$-rays is decreases. & (4) & The cut-off threshold wavelength decreases. \\\hline
\multicolumn{2}{|c|}{} & (5) & The cut-off threshold wavelength remains unchanged. \\\hline
\end{tabular}
\end{center}
\begin{qoptions}
\opt{A}{P $\to$ 3,5; Q $\to$ 1,4; R $\to$ 1; S $\to$ 2}
\opt{B}{P $\to$ 2,4; Q $\to$ 2; R $\to$ 4; S $\to$ 3}
\opt{C}{P $\to$ 3,5; Q $\to$ 3; R $\to$ 2; S $\to$ 1}
\opt{D}{P $\to$ 3,4; Q $\to$ 2; R $\to$ 3; S $\to$ 1}
\end{qoptions}
\qn{10} Select the matching list.
\begin{center}\footnotesize\renewcommand{\arraystretch}{1.3}\setlength{\tabcolsep}{3pt}
\begin{tabular}{|l@{\hspace{3pt}}p{2.4cm}|l@{\hspace{3pt}}p{2.4cm}|}
\hline \multicolumn{2}{|c|}{List-I} & \multicolumn{2}{c|}{List-II} \\\hline
(P) & Convex mirror, virtual object. & (1) & Real image. \\\hline
(Q) & Concave mirror, virtual object. & (2) & Virtual image. \\\hline
(R) & Concave lens, real object. & (3) & Enlarged image \\\hline
(S) & Convex lens, real object. & (4) & Diminished image \\\hline
\multicolumn{2}{|c|}{} & (5) & Positive focal length. \\\hline
\end{tabular}
\end{center}
\begin{qoptions}
\opt{A}{P $\to$ 1,2,3,4,5; Q $\to$ 1,4; R $\to$ 2,4; S $\to$ 1,2,3,4,5}
\opt{B}{P $\to$ 1,2,5; Q $\to$ 1,3,4; R $\to$ 2,4; S $\to$ 1,2,5}
\opt{C}{P $\to$ 1,2,4,5; Q $\to$ 1,3,4; R $\to$ 2,3,4; S $\to$ 1,2,3,4,5}
\opt{D}{P $\to$ 1,2,3,4,5; Q $\to$ 1,3,4; R $\to$ 2,4; S $\to$ 1,2,3,4,5}
\end{qoptions}
\qn{11} Select the matching list.
\begin{center}\footnotesize\renewcommand{\arraystretch}{1.3}\setlength{\tabcolsep}{3pt}
\begin{tabular}{|l@{\hspace{3pt}}p{2.4cm}|l@{\hspace{3pt}}p{2.4cm}|}
\hline \multicolumn{2}{|c|}{List-I} & \multicolumn{2}{c|}{List-II} \\\hline
(P) & Mosley's law & (1) & Characteristic x-ray. \\\hline
(Q) & Change of photon energy into kinetic energy of electrons & (2) & Photoelectric effect. \\\hline
(R) & Electron emission from a material. & (3) & Hydrogen emission spectrum. \\\hline
(S) & Transition between two atomic energy levels. & (4) & $\gamma$-decay \\\hline
\multicolumn{2}{|c|}{} & (5) & $\alpha$-decay (only) \\\hline
\end{tabular}
\end{center}
\begin{qoptions}
\opt{A}{P $\to$ 1; Q $\to$ 3; R $\to$ 5; S $\to$ 2}
\opt{B}{P $\to$ 1; Q $\to$ 3; R $\to$ 4; S $\to$ 1}
\opt{C}{P $\to$ 1; Q $\to$ 2; R $\to$ 4; S $\to$ 2}
\opt{D}{P $\to$ 1; Q $\to$ 2; R $\to$ 2; S $\to$ 1,3}
\end{qoptions}
\sect{SECTION-II}
\qn{1} Virat got a very large number of really thin convex lenses for Diwali. He placed them very closely one behind the other so that the first one had a focal length $f=10$ cm, the second one $2f$, a third $4f$ and so on. He also found a concave mirror with radius of curvature $f$, which he placed right behind this row of lenses. He inserted a luminous object into the focus of the first lens. What distance (in cm) in front of the mirror did its image appear?
\qn{2} A metallic surface is illuminated alternatively with electromagnetic radiations of wavelengths 3000 \AA{} and 6000 \AA. It is observed that the maximum speeds of the photo electrons under these illuminations are in the ratio $3:1$. Calculate the work function of the metal (in eV). Round off to nearest integer.
\qn{3} A hydrogen atom with its electron in the $n$th orbit is ionized by electromagnetic radiation of wavelength 45 nm. If the kinetic energy of the ejected electron is 26.75 eV then the value of $n$ is ($hc=1242$ eV nm).
\qn{4} A parallel narrow beam of light is incident on a glass sphere of radius $R$ and index 1.5 as shown in figure. A plane mirror is at distance 18 m from the centre of sphere. Find the distance of final image from the plane mirror. (Take $R=8$ m)
\qfig{p1_phy_s4_q04_1}
\qn{5} A person looking into an empty container is able to see the far edge of the container's bottom as shown. The height of the container is $h$, and its width is $d$. When the container is completely filled with a fluid of index of refraction $n$ and viewed from the same angle, the person can see the center of a coin at the middle of the container's bottom as shown. Find $(h/d)^2$ if $n=\sqrt3$.
\qfig{p1_phy_s4_q05_1}
\qn{6} For a camera lens, ``depth of field'' is how far a point object can be away from the position where it would be precisely in focus and still have the light from it fall on the film within a ``circle of confusion'' of some diameter, say $\ell$. For a given picture $\Delta q$ is depth of field, the focal length of the lens is $f$ and $D$ is diameter of lens. Find depth of field (in m). (You may consider the object distance to be much larger than the focal length.) ($\ell=1$ mm, $D=5$ cm, $q=5$ m, $f=50$ cm)
\qfig{p1_phy_s4_q06_1}
\partsec{2}{CHEMISTRY}
\sect{SECTION-I(i)}
\qn{1} Morphine changes into Heroin by replacing
\begin{qoptions}
\opt{A}{Both phenolic and alcoholic $-$OH by $-$OCH$_3$}
\opt{B}{Phenolic $-$OH by $-$OCH$_3$}
\opt{C}{Both phenolic and alcoholic $-$OH by $-$OAc}
\opt{D}{Alcoholic $-$OH by $-$OCH$_3$}
\end{qoptions}
\qn{2} \qfig{p1_chem_s1_q01_1}
The name of polymer is :
\begin{qoptions}
\opt{A}{Polystyrene}
\opt{B}{BuNa$-$S}
\opt{C}{BuNa$-$N}
\opt{D}{BuNa$-$CN}
\end{qoptions}
\qn{3} If DNA fragment is $3'$-\texttt{ATGACTGGA}-$5'$ then it's transcribed RNA fragment is :
\begin{qoptions}
\opt{A}{$3'$-\texttt{ATGCTCCAT}-$5'$}
\opt{B}{$3'$-\texttt{UACUGACCU}-$5'$}
\opt{C}{$5'$-\texttt{UACUGACCU}-$3'$}
\opt{D}{$3'$-\texttt{UTCUGUCCU}-$5'$}
\end{qoptions}
\qn{4} Hydrolysis of sucrose into (+) glucose and (-) fructose is known as-
\begin{qoptions}
\opt{A}{Mutarotation}
\opt{B}{Inversion}
\opt{C}{Pyrolysis}
\opt{D}{None of these}
\end{qoptions}
\sect{SECTION-I(ii)}
\qn{5} One mole of the following compound(s) can consume more than 4 moles of HI.
\begin{qoptions}
\opt{A}{\qfigi{p1_chem_s2_q05_A}}
\opt{B}{\qfigi{p1_chem_s2_q05_B}}
\opt{C}{\qfigi{p1_chem_s2_q05_C}}
\opt{D}{\qfigi{p1_chem_s2_q05_D}}
\end{qoptions}
\qn{6} Which reagent can be used to distinguish t-Butyl alcohol and isopropyl alcohol :
\begin{qoptions}
\opt{A}{KMnO$_4$}
\opt{B}{HCl + Anhy ZnCl$_2$}
\opt{C}{Ammonical AgNO$_3$}
\opt{D}{K$_2$Cr$_2$O$_7$}
\end{qoptions}
\qn{7} Following statement(s) is/are true for D-glucose :
\begin{qoptions}
\opt{A}{It shows Schiff's reagent test}
\opt{B}{It is C-2 epimer of D-mannose}
\opt{C}{Oxidise with Br$_2$ / H$_2$O to form saccharic acid}
\opt{D}{Shows mutarotation.}
\end{qoptions}
\end{multicols}
\noindent\begin{minipage}{\textwidth}\sect{SECTION-I(iii)}
\qn{8}
\begin{center}\footnotesize\renewcommand{\arraystretch}{1.3}\setlength{\tabcolsep}{3pt}
\begin{tabular}{|l@{\hspace{3pt}}p{2.4cm}|l@{\hspace{3pt}}p{2.4cm}|}
\hline \multicolumn{2}{|c|}{List-I} & \multicolumn{2}{c|}{List-II} \\\hline
(P) & Ethynylestradiol & (1) & Analgesic \\\hline
(Q) & Salvarsan & (2) & Tranquilisers \\\hline
(R) & Equanil & (3) & Antibiotics \\\hline
(S) & Paracetamol & (4) & Antifertility drug \\\hline
\multicolumn{2}{|c|}{} & (5) & Antipyretic \\\hline
\end{tabular}
\end{center}
\begin{qoptions}
\opt{A}{P $\to$ 1; Q $\to$ 2,5; R $\to$ 3; S $\to$ 4}
\opt{B}{P $\to$ 1; Q $\to$ 4; R $\to$ 2,5; S $\to$ 3}
\opt{C}{P $\to$ 4; Q $\to$ 3; R $\to$ 2; S $\to$ 1,5}
\opt{D}{P $\to$ 2; Q $\to$ 3; R $\to$ 4; S $\to$ 1,5}
\end{qoptions}
\end{minipage}
\par\medskip
\noindent\begin{minipage}{\textwidth}
\qn{9} List-I contains four isomers of C$_4$H$_{10}$O and List-II contains reactions \& properties.
\par\medskip
\begin{center}\small\renewcommand{\arraystretch}{1.3}
\begin{tabular}{|l@{\hspace{4pt}}c|l@{\hspace{4pt}}p{6.4cm}|}
\hline \multicolumn{2}{|c|}{List-I} & \multicolumn{2}{c|}{List-II} \\\hline
(P) & \qfigi{p1_chem_s3_q09_1} & (1) & Produces H$_2$ on reaction with metallic Na \\\hline
(Q) & \qfigi{p1_chem_s3_q09_2} & (2) & Undergoes oxidation to aldehyde or ketone by PCC \\\hline
(R) & \qfigi{p1_chem_s3_q09_3} & (3) & Cloudiness appears immediately in Lucas Test at room temperature \\\hline
(S) & CH$_3$CH$_2$--O--CH$_2$CH$_3$ & (4) & Reaction with conc. HI it gives mixture of ethyl alcohol and ethyl iodide \\\hline
\multicolumn{2}{|c|}{} & (5) & Gives red color in Victor Meyer's test. \\\hline
\end{tabular}
\end{center}
\begin{qoptions}
\opt{A}{P $\to$ 1,2,5; Q $\to$ 1,3; R $\to$ 1,2; S $\to$ 4}
\opt{B}{P $\to$ 1,2,3; Q $\to$ 1,3; R $\to$ 1,3; S $\to$ 5}
\opt{C}{P $\to$ 2,3,4; Q $\to$ 1,4; R $\to$ 2,3; S $\to$ 4,5}
\opt{D}{P $\to$ 1,4; Q $\to$ 2,5; R $\to$ 2,4,5; S $\to$ 1,2}
\end{qoptions}
\end{minipage}
\par\medskip
\noindent\begin{minipage}{\textwidth}
\qn{10} Match the List I with List II and select the correct answer for major product using the codes given below the Lists.
\par\medskip
\begin{center}\small\renewcommand{\arraystretch}{1.3}
\begin{tabular}{|l@{\hspace{4pt}}c|l@{\hspace{4pt}}p{5.2cm}|}
\hline \multicolumn{2}{|c|}{List-I} & \multicolumn{2}{c|}{List-II} \\\hline
(P) & \qfigi{p1_chem_s3_q10_1} & (1) & Elimination Reaction \\\hline
(Q) & \qfigi{p1_chem_s3_q10_2} & (2) & Carbocation intermediate \\\hline
(R) & \qfigi{p1_chem_s3_q10_3} & (3) & Reaction follows I$^{\text{st}}$ order kinetics \\\hline
(S) & \qfigi{p1_chem_s3_q10_4} & (4) & Reaction follows II$^{\text{nd}}$ order kinetics \\\hline
\multicolumn{2}{|c|}{} & (5) & Reaction is solvolysis \\\hline
\end{tabular}
\end{center}
The correct option is
\begin{qoptions}
\opt{A}{P $\to$ 1,2; Q $\to$ 1,2,3; R $\to$ 2,3; S $\to$ 4}
\opt{B}{P $\to$ 1,2,4; Q $\to$ 1,2,3; R $\to$ 2,3; S $\to$ 3,4}
\opt{C}{P $\to$ 1,2; Q $\to$ 1,3; R $\to$ 2,3; S $\to$ 4,5}
\opt{D}{P $\to$ 1,4; Q $\to$ 1,2,3; R $\to$ 2,3,5; S $\to$ 4}
\end{qoptions}
\end{minipage}
\par\medskip
\noindent\begin{minipage}{\textwidth}\qn{11} Match the list
\begin{center}\footnotesize\renewcommand{\arraystretch}{1.3}\setlength{\tabcolsep}{3pt}
\begin{tabular}{|l@{\hspace{3pt}}p{2.2cm}|l@{\hspace{3pt}}p{3.2cm}|}
\hline \multicolumn{2}{|c|}{List-I} & \multicolumn{2}{c|}{List-II} \\\hline
(P) & Acetophenone & (1) & Reaction with AgNO$_3$ + NH$_4$OH to give white ppt. \\\hline
(Q) & Benzaldehyde & (2) & Reaction with NaOI to give yellow ppt. of CHI$_3$ \\\hline
(R) & Butanal & (3) & Reaction with Fehling to give brown ppt. of Cu$_2$O \\\hline
(S) & 1-butyne & (4) & Reaction with Ac$_2$O + AcONa to give Cinnamic acid \\\hline
\multicolumn{2}{|c|}{} & (5) & Reaction with Tollen's reagent give silver mirror \\\hline
\end{tabular}
\end{center}
The correct code for your answer :-
\begin{qoptions}
\opt{A}{P $\to$ 2; Q $\to$ 4; R $\to$ 3; S $\to$ 1}
\opt{B}{P $\to$ 1; Q $\to$ 4; R $\to$ 3; S $\to$ 2}
\opt{C}{P $\to$ 2; Q $\to$ 3; R $\to$ 4; S $\to$ 1}
\opt{D}{P $\to$ 1,4; Q $\to$ 2; R $\to$ 4; S $\to$ 3}
\end{qoptions}
\end{minipage}
\qfullstart
\sect{SECTION-II}
\qn{1} \qfig{p1_chem_s4_q01_1}
\qfullsplit
How many moles of HIO$_4$ will be consumed by one mole of compound (A)?
\qfullend
\qn{2} \qfig{p1_chem_s4_q01_2}
The degree of unsaturation of compound (B) is \underline{\hspace{1.5cm}}
\qfullbegin
\qn{3} How many total natural and non essential Amino acids are obtain after complete acidic hydrolysis of following compound.
\qfig{p1_chem_s4_q03_1}
\qfullend
\qn{4} The alcohol among following which show dehydration reaction faster than 2-methyl propan-2-ol ((CH$_3$)$_3$C--OH).
\qfig{p1_chem_s4_q04_1}
\qfig{p1_chem_s4_q04_2}
\noindent\begin{minipage}{\columnwidth}\qn{5} Identify compounds give positive test with NaHSO$_3$?
\par\smallskip
\begin{center}\qfigi{p1_chem_s4_q05_1}\par\smallskip\qfigi{p1_chem_s4_q05_2}\par\smallskip\qfigi{p1_chem_s4_q05_3}\end{center}
\end{minipage}
\par\smallskip
\noindent\begin{minipage}{\columnwidth}\qn{6} Identify compounds give Biurate test?
\par\smallskip
(i) Cellulose \quad (ii) Protien \quad (iii) Nylon-6 \quad (iv) Starch
\par\smallskip
\begin{center}\qfigi{p1_chem_s4_q06_1}\end{center}
\end{minipage}
\partsec{3}{MATHEMATICS}
\sect{SECTION-I(i)}
\qn{1} If $f(x)$ is a polynomial function and $(f(\alpha))^2+(f'(\alpha))^2=0$ then find $\displaystyle\lim_{x\to\alpha}\dfrac{f(x)}{f'(x)}\left[\dfrac{f'(x)}{f(x)}\right]$ (where $[\,.\,]$ denotes greatest integer function)
\begin{qoptions}
\opt{A}{0}
\opt{B}{1}
\opt{C}{$-1$}
\opt{D}{3}
\end{qoptions}
\qn{2} The function $f(x)=\sin^{-1}\left(\dfrac{2x}{1+x^2}\right)$ is
\begin{qoptions}
\opt{A}{Continuous but not differentiable at $x=1$}
\opt{B}{Differentiable at $x=1$}
\opt{C}{Neither continuous nor differentiable at $x=1$}
\opt{D}{Differentiable everywhere}
\end{qoptions}
\qn{3} Range of the function $f(x)=\dfrac{2x^4-14x^2+8x+49}{x^4-7x^2-4x+23}$ is
\begin{qoptions}
\opt{A}{$[1,3]$}
\opt{B}{$[2,3]$}
\opt{C}{$(2,3]$}
\opt{D}{$(1,2]$}
\end{qoptions}
\qn{4} If
\begin{center}\resizebox{\columnwidth}{!}{$\displaystyle f(\theta)=\dfrac{(2\cos\theta-1)(2\cos2\theta-1)(2\cos4\theta-1)\ldots(2\cos2^{n-1}\theta-1)}{2\cos2^n\theta+1}$}\end{center}
for $n\in N$ and $\theta\ne2m\pi\pm\dfrac{2\pi}{3}$, $m\in I$ then $f\left(\dfrac{\pi}{4}\right)$ is equal to
\begin{qoptions}
\opt{A}{$\sqrt2-1$}
\opt{B}{$1-\sqrt2$}
\opt{C}{$\sqrt3-1$}
\opt{D}{$1-\sqrt3$}
\end{qoptions}
\sect{SECTION-I(ii)}
\qn{5} Suppose that $f(a)>0$ and that $f$ is differentiable at $x=a$. Which of the following statements is/are TRUE.
\begin{qoptions}
\opt{A}{$\displaystyle\lim_{n\to\infty}\left(\dfrac{f\left(a+\frac1n\right)}{f(a)}\right)^{\frac1n}=1$}
\opt{B}{$\displaystyle\lim_{n\to\infty}\left(\dfrac{f\left(a+\frac1n\right)}{f(a)}\right)^{\frac1n}=e^{\frac{f'(a)}{f(a)}}$}
\opt{C}{$\displaystyle\lim_{x\to a^+}\left(\dfrac{f(x)}{f(a)}\right)^{\frac{1}{\ln x-\ln a}}=e^{\frac{af'(a)}{f(a)}};\ a>0$}
\opt{D}{$\displaystyle\lim_{x\to a^-}\left(\dfrac{f(x)}{f(a)}\right)^{\frac{1}{\ln x-\ln a}}=e^{\frac{af'(a)}{f(a)}};\ a>0$}
\end{qoptions}
\qn{6} Which of the following equations have no real solutions ?
\begin{qoptions}
\opt{A}{$x^2-2x+5+\pi^x=0$}
\opt{B}{$\log_{1.5}(\cot^{-1}x+\operatorname{sgn}(e^x))=0$}
\opt{C}{$x^4-2x^2\sin^2\dfrac{\pi x}{2}+1=0$}
\opt{D}{$\tan\left(x+\dfrac{\pi}{6}\right)=2\tan x$}
\end{qoptions}
\qn{7} Function $f(x)$ satisfies the relation $f(x)+2f(1-x)=x^2+1\ \forall\ x\in R$ then :
\begin{qoptions}
\opt{A}{$f(x)$ is an odd function}
\opt{B}{$f(x)$ is many-one function}
\opt{C}{Range of $f(x)$ is $\left[-\dfrac13,\infty\right)$}
\opt{D}{$f(x)$ is increasing in the interval $(2,\infty)$}
\end{qoptions}
\end{multicols}
\noindent\begin{minipage}{\textwidth}\sect{SECTION-I(iii)}
\qn{8} If $f(x)$ is twice differentiable function in $[c_1-1,\,c_2+1]$ and $f'(c_1)=f'(c_2)=0$, $f''(c_1).f''(c_2)<0$, $f(c_1)=9$, $f(c_2)=0$. Let $k$ and $m$ be the minimum number of the roots of $f(x)=0$ and $f'(x)=0$ in $[c_1-1,\,c_2+1]$.
\begin{center}\footnotesize\renewcommand{\arraystretch}{1.3}\setlength{\tabcolsep}{3pt}
\begin{tabular}{|l@{\hspace{3pt}}p{3.6cm}|l@{\hspace{3pt}}l|}
\hline \multicolumn{2}{|c|}{List-I} & \multicolumn{2}{c|}{List-II} \\\hline
(P) & If $f''(c_1)-f''(c_2)>0$, then $k=$ & (1) & 1 \\\hline
(Q) & If $f''(c_1)-f''(c_2)<0$, then $k=$ & (2) & 2 \\\hline
(R) & If $f''(c_1)-f''(c_2)>0$, then $m=$ & (3) & 3 \\\hline
(S) & If $f''(c_1)-f''(c_2)<0$, then $m=$ & (4) & 4 \\\hline
\multicolumn{2}{|c|}{} & (5) & 7 \\\hline
\end{tabular}
\end{center}
\begin{qoptions}
\opt{A}{P $\to$ 1; Q $\to$ 4; R $\to$ 2; S $\to$ 3}
\opt{B}{P $\to$ 4; Q $\to$ 1; R $\to$ 3; S $\to$ 1}
\opt{C}{P $\to$ 2; Q $\to$ 1; R $\to$ 4; S $\to$ 2}
\opt{D}{P $\to$ 1; Q $\to$ 4; R $\to$ 3; S $\to$ 4}
\end{qoptions}
\end{minipage}
\par\medskip
\noindent\begin{minipage}{\textwidth}\qn{9} $\alpha,\beta$ are real and distinct roots of $x^2+4bx+c=0$ where $b,c\in R$.
\begin{center}\footnotesize\renewcommand{\arraystretch}{1.3}\setlength{\tabcolsep}{3pt}
\begin{tabular}{|l@{\hspace{3pt}}p{2.9cm}|l@{\hspace{3pt}}p{2.7cm}|}
\hline \multicolumn{2}{|c|}{List-I} & \multicolumn{2}{c|}{List-II} \\\hline
(P) & $f(x)=x^3+6bx^2+3cx+2$; $b<0$, $c>0$, $f(\alpha)f(\beta)>0$ & (1) & $f(x)=0$ has three real roots out of which two are positive real roots and one is negative real root \\\hline
(Q) & $f(x)=cx^3+6bx^2+3x-2$; $c<0$; $f(1/\alpha)f(1/\beta)>0$ & (2) & $f(x)=0$ has exactly one real root which is positive \\\hline
(R) & $f(x)=x^3+6x^2+3cx-2$; $c<0$, $f(\alpha)f(\beta)>0$ & (3) & $f(x)=0$ has three real roots out of which two are negative real roots and one is positive real root \\\hline
(S) & $f(x)=cx^3+6bx^2+3x-2$; $c<0$, $f(1/\alpha)f(1/\beta)<0$ & (4) & $f(x)=0$ has exactly one real root which is negative \\\hline
\multicolumn{2}{|c|}{} & (5) & $f(x)=0$ has three positive real roots \\\hline
\end{tabular}
\end{center}
\begin{qoptions}
\opt{A}{P $\to$ 4; Q $\to$ 3; R $\to$ 1; S $\to$ 3}
\opt{B}{P $\to$ 2; Q $\to$ 4; R $\to$ 1; S $\to$ 3}
\opt{C}{P $\to$ 4; Q $\to$ 4; R $\to$ 2; S $\to$ 1}
\opt{D}{P $\to$ 2; Q $\to$ 3; R $\to$ 2; S $\to$ 1}
\end{qoptions}
\end{minipage}
\par\medskip
\noindent\begin{minipage}{\textwidth}
\qn{10} Given $f_1(x)=x^2\operatorname{sgn}(x)=\begin{cases}x^2 & x\ge0\\ -x^2 & x<0\end{cases}$ ; $f_2(x)=\displaystyle\int_0^x t^2\sin\left(\dfrac1t\right)dt$; $f_2(0)=0$ ; $f_3(x)=x^{-\frac13}\,|\sin x|$, $f_3(0)=0$ ; $f_4(x)=\begin{cases}0 & -1<x\le0\\ -x^3 & 0<x<1\end{cases}$.
\par\medskip
\begin{center}\small\renewcommand{\arraystretch}{1.3}
\begin{tabular}{|l@{\hspace{4pt}}l|l@{\hspace{4pt}}p{6.4cm}|}
\hline \multicolumn{2}{|c|}{List-I} & \multicolumn{2}{c|}{List-II} \\\hline
(P) & The function $f_1$ is & (1) & continuous but not differentiable at $x=0$ \\\hline
(Q) & The function $f_2$ is & (2) & first derivative exists at $x=0$ but second derivative does not exist \\\hline
(R) & The function $f_3$ is & (3) & second derivative exists at $x=0$, but it is not continuous there \\\hline
(S) & The function $f_4$ is & (4) & second derivative exists at $x=0$ and is continuous \\\hline
\multicolumn{2}{|c|}{} & (5) & Discontinuous at $x=0$ \\\hline
\end{tabular}
\end{center}
The correct option is :
\begin{qoptions}
\opt{A}{P $\to$ 2; Q $\to$ 3; R $\to$ 4; S $\to$ 1}
\opt{B}{P $\to$ 2; Q $\to$ 3; R $\to$ 1; S $\to$ 4}
\opt{C}{P $\to$ 1; Q $\to$ 3; R $\to$ 4; S $\to$ 2}
\opt{D}{P $\to$ 1; Q $\to$ 3; R $\to$ 2; S $\to$ 4}
\end{qoptions}
\end{minipage}
\par\medskip
\noindent\begin{minipage}{\textwidth}
\qn{11} If $f(x)=\sin^{-1}\left(\dfrac{4x}{4+x^2}\right)$, then
\par\medskip
\begin{center}\small\renewcommand{\arraystretch}{1.5}
\begin{tabular}{|l@{\hspace{4pt}}l|l@{\hspace{4pt}}l|}
\hline \multicolumn{2}{|c|}{List-I} & \multicolumn{2}{c|}{List-II} \\\hline
(P) & $f\left(\sqrt8-2\right)$ & (1) & $\pi+\tan^{-1}\left(\dfrac{4\pi}{4-\pi^2}\right)$ \\\hline
(Q) & $f\left(-\dfrac{3}{\sqrt8-2}\right)$ & (2) & $-\pi+2\tan^{-1}\left(\dfrac34\left(\sqrt2+1\right)\right)$ \\\hline
(R) & $f\left(\sqrt{12}-\sqrt3\right)$ & (3) & $\tan^{-1}\dfrac{4\pi}{\pi^2-4}$ \\\hline
(S) & $f(\pi)$ & (4) & $\dfrac\pi2-2\tan^{-1}\left(\sqrt2-1\right)$ \\\hline
\multicolumn{2}{|c|}{} & (5) & $\dfrac\pi2-\tan^{-1}\left(\dfrac{\sqrt3}{12}\right)$ \\\hline
\end{tabular}
\end{center}
\begin{qoptions}
\opt{A}{P $\to$ 4; Q $\to$ 2; R $\to$ 5; S $\to$ 3}
\opt{B}{P $\to$ 5; Q $\to$ 4; R $\to$ 2; S $\to$ 1}
\opt{C}{P $\to$ 2; Q $\to$ 4; R $\to$ 5; S $\to$ 1}
\opt{D}{P $\to$ 1; Q $\to$ 2; R $\to$ 4; S $\to$ 3}
\end{qoptions}
\end{minipage}
\begin{multicols}{2}
\sect{SECTION-II}
\qn{1} Find the number of real solution of the equation $x^3+3x^2-(x-6)^{1/3}+8+3x=0$
\qn{2} The number of solutions of the equation $\left(x-2.\left[\dfrac x2\right]-1\right)^2=\dfrac12$ for $x\in[-4,4]$, is (where $[\,.\,]$ denotes greatest integer function)
\qn{3} Let $f:A\to\left[-\dfrac12,0\right]$ is such that $f^3(x)-3xf(x)+x^3+1=0$ and $f(x)$ is an invertible function, then number of solution(s) of the equation $(f(f(x)))^2=\dfrac32f^{-1}(x)+1$ is
\qn{4} Let $P(x)$ be a polynomial of degree 4 having a local maximum at $x=2$ and $\displaystyle\lim_{x\to0}\left(3-\dfrac{P(x)}{x}\right)=27$. If $P(1)=-9$ and $P''(x)$ has a local minimum at $x=2$. If global maximum value of $y=P'(x)$ on the set $A=\{x:x^2+12\le7x\}$ is $4M$ then value of $M$ is :
\qn{5} Number of solution of the equation $\left(\cos^{-1}\left(\dfrac{1-4x^2}{1+4x^2}\right)\right)^{\frac34}=x^2-x\sin x-\cos x$ is
\qn{6} Let $f(x)$ is defined only for $x\in(0,5)$ and defined as $f^2(x)=1\ \forall\ x\in(0,5)$. Function $f(x)$ is continuous for all $x\in(0,5)-\{1,2,3,4\}$ (at $x=1,2,3,4$ $f(x)$ may or may not be continuous). The number of possible function $f(x)$ if it is discontinuous at two and three integral points in $(0,5)$ are $a$ and $b$ respectively, then $a+b$ is equal to
\papersec{2}
\partsec{1}{PHYSICS}
\sect{SECTION-I(i)}
\qn{1} The half-life of radioactive sample is 1620 years and its atomic mass is 226 The number of atoms that will decay from its 1 gm sample per second will be
\begin{qoptions}
\opt{A}{$3.61\times10^{10}$}
\opt{B}{$3.6\times10^{12}$}
\opt{C}{$3.11\times10^{15}$}
\opt{D}{$31.1\times10^{15}$}
\end{qoptions}
\qn{2} Kuldeep has a well and he constantly draws water from it with a volumetric flow rate $q$. Water is being constantly replenished into the well so the total volume stays constant at $V$. However, Kuldeep accidentally dropped a piece of radioactive bread into the well, which dissipated perfectly in the water. The bread contained radioactive isotopes with a half-life $T$. Kuldeep decided to ignore it and continued to draw water at the original rate. How long does it take until the radioactive activity of the well drops below $A$? ($N_0$ are the initial no. of nuclei.)
\begin{qoptions}
\opt{A}{$\dfrac{V}{q}\ln\left(\dfrac{N_0\ln2}{AT}\right)$}
\opt{B}{$\dfrac{1}{\frac{\ln2}{T}-\frac{q}{V}}\ln\left(\dfrac{N_0\ln4}{AT}\right)$}
\opt{C}{$\dfrac{1}{\frac{\ln2}{T}+\frac{q}{V}}\ln\left(\dfrac{N_0\ln2}{AT}\right)$}
\opt{D}{None}
\end{qoptions}
\qn{3} Far from a material with a thickness $t$, we place a source of gamma rays with an energy of approximately 2 MeV. The radiation has an absorption coefficient of $\mu_1$ in the material. For radiation of this energy can happen that a photon converts into an electron-positron pair inside the material (the loss in intensity due to this effect is already included in the $\mu_1$ coefficient). However, the positron immediately annihilates with an electron and produces 2 gamma flashes, each with energy $E=511$ keV. The absorption coefficient of the material for the radiation with this energy is $\mu_2$. We place a detector far behind the material and observe the intensity of radiation of the energy $E$. For what thickness of material $t$ will this intensity be the largest? Given that: $I=I_0e^{-\mu x}$ where $I_0$ is the intensity before entering the material and $x$ is the depth at which the intensity is calculated.
\begin{qoptions}
\opt{A}{$\dfrac{3}{\mu_2-\mu_1}\ln\left(\dfrac{\mu_2}{\mu_1}\right)$}
\opt{B}{$\dfrac{3}{\mu_2-\mu_1}\ln\left(\dfrac{2\mu_2}{\mu_1}\right)$}
\opt{C}{$\dfrac{1}{\mu_2-\mu_1}\ln\left(\dfrac{\mu_2}{\mu_1}\right)$}
\opt{D}{None}
\end{qoptions}
\qn{4} Binding energy per nucleon versus mass number curve for nuclei is shown in figure. A, B, C and D are four nuclei indicated on the curve. The $Q$ value of following nuclear reaction $A+B\to C$ is $N\times10^7$ eV. Find the value of $N$.
\qfig{p2_phy_s1_q04_1}
\begin{qoptions}
\opt{A}{210}
\opt{B}{21}
\opt{C}{2.1}
\opt{D}{None}
\end{qoptions}
\sect{SECTION-I(ii)}
\qn{5} Object-A is at a distance $2f$ from a convex lens of focal length $f$, and Object-B is at $f/2$ from the lens. On the other side of the lens there is a flat mirror laying parallel to and coincide with the optical axis of the lens. Consider refraction from lens and then reflection from plane mirror.
\qfig{p2_phy_s2_q05_1}
\begin{qoptions}
\opt{A}{The images are formed at a distance of $2f$ and $f$ from lens.}
\opt{B}{The magnitude of magnification is 1 in both cases}
\opt{C}{The magnitude of magnifications are 1 and 2.}
\opt{D}{Both the images are virtual.}
\end{qoptions}
\qn{6} We can imagine the electron (charge: $e$) as a thin spherical shell with uniformly distributed charge and radius $R$. For the electron, a magnetic moment can be created by making the sphere rotate around an axis passing through its centre. The electron's magnetic moment is known to be $\mu=\dfrac{eh}{4\pi m}$, where $h$ is the Planck constant, ($c$: speed of light)
\begin{qoptions}
\opt{A}{The maximum magnetic moment that the sphere can have is $\dfrac13ecR$.}
\opt{B}{The minimum possible radius of the electron is $\dfrac{3h}{4\pi mc}$.}
\opt{C}{Maximum energy that can be stored in the electron is $\dfrac{ke^2}{4R}$.}
\opt{D}{none}
\end{qoptions}
\qn{7} Consider three converging lenses $L_1$, $L_2$ and $L_3$ having identical geometrical construction. The index of refraction of $L_1$ and $L_2$ are $\mu_1$ and $\mu_2$ respectively. The upper half of the lens $L_3$ has a refractive index $\mu_1$ and the lower half has $\mu_2$ as shown in the figure. A point object $O$ is imaged at $O_1$ by the lens $L_1$ and at $O_2$ by the lens $L_2$ placed in same position. If $L_3$ is placed at the same place,
\qfig{p2_phy_s2_q07_1}
\begin{qoptions}
\opt{A}{There will be an image at $O_1$.}
\opt{B}{There will be an image at $O_2$.}
\opt{C}{The only image will form somewhere between $O_1$ and $O_2$.}
\opt{D}{The only image will form away from $O_2$.}
\end{qoptions}
\sect{SECTION-II(i)}
\paraq{1}{2}
A thin equiconvex lens of refractive index $\dfrac32$ is placed on a horizontal plane mirror as shown in the figure. The space between lens and the mirror is then filled with water of refractive index $\dfrac43$. It is found that when a point object is placed 15 cm above the lens on the principal axis the object coincides with its own image.
\qfig{p2_phy_s3_para1_2_1}
\qn{1} At what distance $xR$ object can be placed before water is filled so that image coincides with object if $R$ is radius of curvature of each surface of lens. Find $x$.
\qn{2} What can be the value of $R$ (in cm) from above information?
\paraq{3}{4}
Two identical objects $AB$ and $CD$ each of height 1 mm are placed at separation of 80 cm. A converging lens of focal length of 30 cm with its principal axis along line $AC$ is placed such that virtual image of point $A$ and real image of point $C$ coincide.
\qfig{p2_phy_s3_para3_4_1}
\qn{3} Distance (in cm) of point $A$ from lens is :
\qn{4} If image of point $B$ is $B'$ and $D$ is $D'$ then length of $B'D'$ (in mm) is :
\sect{SECTION-II(ii)}
\qn{5} Consider two lenses at a mutual distance $d=50$ cm. The first is a biconvex lens with focal length 20 cm, and the second is a biconcave lens with focal length $-20$ cm. Approximately how far (in cm) from the first lens do we need to place the object for the system to project it on itself?
\qn{6} A right-angled isosceles prism of side length 9cm and refractive index 1.5 is placed 6cm away from a converging lens of focal length $f_1=20$ cm, followed by a diverging lens of focal length $f_2=-10$ cm a distance 7cm behind it. A 1cm stick is located 8cm above the prism, with one end aligned with the mid-point of the hypotenuse as shown in the figure below. The final image is formed at a distance $x$ from the concave lens. Find $x$ (in cm):
\qfig{p2_phy_s4_q06_1}
\qn{7} A pont source of power 0.1W and emitting light of wavelength 5000 \AA{} is placed at a distance 20 cm from a sphere of radius 30 cm as shown in the figure. The work function of the sphere is 1.5 eV. The efficiency of photoelectric emission of the sphere is $10^{-6}$. If the time from the beginning of the experiment after which the photoelectric current stops is approximately $m\times10^n$ sec in scientific notation. ($m$ and $n$ are integer). Find $|m|+|n|$
\qfig{p2_phy_s4_q07_1}
\qn{8} Chlorines molecule, which is moving with a velocity $v=600$ m/s, absorbs a photon with a wavelength $\lambda=350$ nm and is divided into two atoms. The first atom's velocity is measured to be $u=1600$ m/s, which perpendicular to the molecule's initial velocity. Find the chlorine molecule's binding energy if all the particles were at a state with minimal inner energy. It is presumed that a photon's momentum is negligible compared to chlorine's momentum. If Chlorine molecule's binding energy is approximately $x$ eV then report $2x$ as answer.
\qn{9} For kinetic energy $K_0$ a free proton just is a able to ionize a hydrogen atom. Assume that the electron inside the hydrogen atom is still and that the electromagnetic interaction between the atomic nucleus and the free proton is negligible. Hydrogen's binding energy is $E_0=13.6$ eV, protons mass is $m_p=1.67\times10^{-27}$ kg and electrons mass $m_e=9.11\times10^{-31}$ kg. If the smallest kinetic energy is $K_0$ keV then report $4K_0$ as answer. Round off to nearest integer.
\qn{10} A light ray starts from the interior of the circumference of a solid circle of radius $R$ at $\pi/4$ radians with respect to the radial direction, radially inwards. Set the origin to be at the center of the circle. If the refractive index of the circle varies according to the relationship $n(r)=\sqrt{\dfrac{R^2}{r^2}+1}$, determine the magnitude of the angular displacement (in rad) of the light ray when it reaches the center of the circle. ($r$ is the distance from the centre of circle)
\partsec{2}{CHEMISTRY}
\sect{SECTION-I(i)}
\qn{1} \qfig{p2_chem_s1_q01_1}
\begin{qoptions}
\opt{A}{Anti-aromatic}
\opt{B}{Non-aromatic}
\opt{C}{Aromatic}
\opt{D}{Quasi-aromatic}
\end{qoptions}
\qn{2} Correct reactivity order for S$_\text{N}$1 reaction
\qfig{p2_chem_s1_q02_1}
\begin{qoptions}
\opt{A}{I $>$ II $>$ IV $>$ III}
\opt{B}{II $>$ IV $>$ III $>$ I}
\opt{C}{I $>$ II $>$ III $>$ IV}
\opt{D}{II $>$ IV $>$ I $>$ III}
\end{qoptions}
\qn{3} Which reaction shows correct product.
\begin{qoptions}
\opt{A}{\qfigi{p2_chem_s1_q03_A}}
\opt{B}{\qfigi{p2_chem_s1_q03_B}}
\opt{C}{\qfigi{p2_chem_s1_q03_C}}
\opt{D}{CH$_3$--CH$_2$--OH $\xrightarrow{\text{PCC}}$ CH$_3$ -- COOH}
\end{qoptions}
\qn{4} The number of isomeric tetraenes (NOT containing sp-hybridized carbon atoms) that can be formed from the following reaction sequence is \underline{\hspace{1.5cm}}
\qfig{p2_chem_s1_q04_1}
\begin{qoptions}
\opt{A}{2}
\opt{B}{4}
\opt{C}{3}
\opt{D}{1}
\end{qoptions}
\sect{SECTION-I(ii)}
\qn{5} Choose the incorrect option(s)
\begin{qoptions}
\opt{A}{Rate of S$_\text{N}$2 \qfigi{p2_chem_s2_q05_B}}
\opt{B}{Ph--CH(OH)--CH(D)--Me $\xrightarrow[\text{S}_\text{N}1]{\text{HI}}$ racemic mixture (Optically pure)}
\opt{C}{Me--Cl $\xrightarrow[\text{Ag}_2\text{O}]{\text{Moist}}$ Me--O--Me (Major)}
\opt{D}{MeI (excess) + NH$_3\to$ Me$_3$N}
\end{qoptions}
\qn{6} Which of the following reaction(s) can generate a five membered ring as the major product?
\begin{qoptions}
\opt{A}{\qfigi{p2_chem_s2_q06_A}}
\opt{B}{\qfigi{p2_chem_s2_q06_B}}
\opt{C}{\qfigi{p2_chem_s2_q06_C}}
\opt{D}{\qfigi{p2_chem_s2_q06_D}}
\end{qoptions}
\qfullbegin
\qn{7} Some organic compounds are given below
\qfig{p2_chem_s2_q07_1}
\qfullsplit
Select the correct statement about the above compounds from the following
\begin{qoptions}
\opt{A}{Compounds I and III would reduce Fehiling's solution and will also mutarotate.}
\opt{B}{Compound IV will not reduce Tollen's reagent and will not mutarotate.}
\opt{C}{Compound II will not be reducing but mutarotating.}
\opt{D}{The product(s) of the acid catalysed hydrolysis of both II and III will be reducing but not mutarotating.}
\end{qoptions}
\qfullend
\sect{SECTION-II(i)}
\paraq{1}{2}
Carbohydrates are essential for life in both plants and animals. They form a major portion of our food. Honey has been used for a long time as an instant source of energy by `Vaids' in ayurvedic system of medicine. Carbohydrates are used as storage molecules as starch in plants and glycogen in animals. Cell wall of bacteria and plants is made up of cellulose.
\qn{1} When the following aldohexose exists in its D-configuration, the total number of stereoisomers in its pyranose form is
\qfig{p2_chem_s3_q01_1}
\qn{2} Consider the following reactions :

{\small (i) Glucose + ROH $\xrightarrow{\text{dry HCl}}$ Acetal $\xrightarrow[\text{(CH}_3\text{CO)}_2\text{O}]{x\text{ eq. of}}$ acetyl derivative

(ii) Glucose $\xrightarrow{\text{Ni/H}_2}$ A $\xrightarrow[\text{(CH}_3\text{CO)}_2\text{O}]{y\text{ eq. of}}$ acetyl derivative

(iii) Glucose $\xrightarrow[\text{(CH}_3\text{CO)}_2\text{O}]{z\text{ eq. of}}$ acetyl derivative}

Report $x+y+z$?
\paraq{3}{4}
In halogenation of an alkane, the alkane is said to undergo fluorination, chlorination, bromination or iodination depending on the identity of the halogen reactant. Chlorination and bromination are the two widely used alkane halogenation reactions. Fluorination reactions generally proceed too quickly to be useful and iodination reactions go too slowly.

Halogenations usually result in the formation of a mixture of products rather than a single product. More than one product results because more than one hydrogen atom on an alkane can be replaced with halogen atoms.
\qn{3} The maximum number of isomers (including stereoisomers) that are possible on -mono-chlorination of the following compounds, is :
\qfig{p2_chem_s3_q03_1}
\qn{4} The maximum number of possible isomers (including stereoisomers) which may be formed on mono-bromination of 1-methylcyclohex-1-ene using NBS is \underline{\hspace{1cm}}
\qfullbegin
\sect{SECTION-II(ii)}
\qn{5} \qfig{p2_chem_s4_q05_1}
\qfullsplit
What is the value of $x+y+z=\ ?$
\qfullend
\qn{6} In the following monobromination reaction, the number of possible chiral products is :
\qfig{p2_chem_s4_q06_1}
\qn{7} The total number of alkenes possible by dehydrobromination of 3-bromo-3-cyclopentylhexane using alcoholic KOH is
\qn{8} How many organic compound(s) is/are correctly matched with number of chiral carbon(s) present in it's structure :

(i) Histamine = 0 \quad (ii) Chloramphenicol = 2

(iii) Penicillin = 3 \quad (iv) Aspartame = 4

(v) Sucralose = 9
\qn{9} A decapeptide (Mol. Wt. 796) on complete hydrolysis gives glycine (Mol. Wt. 75), alanine and phenylalanine. Glycine contributes 47.0\% to the total weight of the hydrolysed products. The number of glycine units present in the decapeptide is
\qn{10} Degree of unsaturation in Aspirin is \underline{\hspace{2.5cm}}
\partsec{3}{MATHEMATICS}
\sect{SECTION-I(i)}
\qn{1} If $[\,.\,]$ denotes the greatest integer function then the domain of the real valued function $f(x)=\log_{\left[x+\frac12\right]}\left|x^2-x-2\right|$ is
\begin{qoptions}
\opt{A}{$\left[\dfrac32,\infty\right)$}
\opt{B}{$\left[\dfrac32,2\right)\cup(2,\infty)$}
\opt{C}{$\left(\dfrac12,2\right)\cup(2,\infty)$}
\opt{D}{$(2,\infty)$}
\end{qoptions}
\qn{2} Number of solution of the equation $|\sin^{-1}(\sin x)|=\cos x$, for $x\in\left[\dfrac{-\pi}{2},\dfrac\pi2\right]$ is equal to
\begin{qoptions}
\opt{A}{0}
\opt{B}{2}
\opt{C}{4}
\opt{D}{6}
\end{qoptions}
\qn{3} Let a line $\ell$ through origin is tangent to the curve $y=x^3+x+16$. The slope of the line $\ell$ is
\begin{qoptions}
\opt{A}{2}
\opt{B}{8}
\opt{C}{13}
\opt{D}{26}
\end{qoptions}
\qn{4} If $\displaystyle\lim_{x\to0}\dfrac{\tan^3x+\tan^32x+\tan^33x}{\tan^34x+\tan^35x+\tan^36x}=\dfrac pq$ ($p,q\in N$), then least value of $(p+q)$ is equal to
\begin{qoptions}
\opt{A}{45}
\opt{B}{47}
\opt{C}{49}
\opt{D}{51}
\end{qoptions}
\sect{SECTION-I(ii)}
\qn{5} If $f(x)=\begin{cases}x.\sin\left(\dfrac\pi x\right) & ,\ x\ne0\\ k & ,\ x=0\end{cases}$ then
\begin{qoptions}
\opt{A}{for $k=1$, $f(x)$ is discontinuous at $x=0$}
\opt{B}{for $k=1$, $f(x)$ is continuous at $x=0$}
\opt{C}{for $k=\pi$, $f(x)$ is continuous at $x=0$}
\opt{D}{for $k=0$, $f(x)$ is continuous at $x=0$}
\end{qoptions}
\qn{6} For a certain curve $y=f(x)$ satisfying $\dfrac{d^2y}{dx^2}=6x-4$, $f(x)$ has local minimum value 5 when $x=1$ then which of the following is/are correct ?
\begin{qoptions}
\opt{A}{Equation of curve $y$ is $y=x^3-2x^2+x+5$}
\opt{B}{$f(x)$ has a local maximum at $x=\dfrac13$}
\opt{C}{Equation of curve is $y=x^3-2x^2-x+7$}
\opt{D}{The number of real and distinct roots of $f(x)=0$ is 1}
\end{qoptions}
\qn{7} The value of `$x$' satisfying
\[\sin(2\cos^{-1}(\cot(\tan^{-1}x)))=0\]
is/are
\begin{qoptions}
\opt{A}{1}
\opt{B}{0}
\opt{C}{$-1$}
\opt{D}{$\dfrac12$}
\end{qoptions}
\sect{SECTION-II(i)}
\paraq{1}{2}
Consider the function $f(x)=x|x^2-x|+|x-2|$.
\qn{1} The total number of local maxima and local minima of $f(x)$ is
\qn{2} If $x=\alpha$ is the point of maxima of the function $f(x)$, then $|6\alpha|=$
\paraq{3}{4}
Let $f(x)=x^2e^x-2e^x+3$, $x\in(-4,1)$
\qn{3} Number of real roots of the equation $f(x)=0$ is
\qn{4} Number of local maxima of the function $f(x)$ is
\sect{SECTION-II(ii)}
\qn{5} `$p$' is number of points where $y=(x^3+1)\,|x^2-1|+\sin|x+1|$ is not differentiable `$q$' is number of points where $y=\operatorname{sgn}(x^3+x^2-x+3)$ is discontinuous and $r=\sqrt3\tan\left(\cos^{-1}\left(\dfrac12\right)+\tan^{-1}\left(\dfrac{-1}{\sqrt3}\right)\right)$ then $p+q+r$ is equal to
\qn{6} $f(x)=\begin{cases}5 & x=0\\ x^2+5x+\alpha & 0<x<2\\ \beta x+\gamma & 2\le x\le3\end{cases}$
\par
If $f(x)$ satisfies conditions of LMVT in $[0,3]$ then $2\alpha-\beta+\gamma+f''(1)+\dfrac{f'(3)}{9}$ is :
\qn{7}
\begin{center}\resizebox{\columnwidth}{!}{$\displaystyle\lim_{x\to0}\left(\dfrac{1-\cos x\sqrt{\cos2x}}{x^2}+\sqrt[3]{\dfrac{1+\cos^2x+3\cos^4x+3\cos x}{\cos x+63}}\right)$}\end{center}
is equal to
\qn{8} Let $f(x)=\tan^{-1}(\cot x-2\cot2x)$ and $\displaystyle\sum_{r=1}^{5}f(r)=a-b\pi$, where $a,b\in N$. Find the value of $(a+b)$
\qn{9} For any real number $x$, let $[x]$ denote the largest integer less than or equal to $x$. Let $f$ be a real valued function defined on the interval $[-3,3]$ by $f(x)=\begin{cases}-x-[-x] & \text{if }[x]\text{ is even}\\ x-[x] & \text{if }[x]\text{ is odd}\end{cases}$.
If $L$ denotes the number of points of discontinuity and $M$ denotes the number of points of non-derivability of $f(x)$, then find $(L+M)$.
\qn{10} If $p_1$ and $p_2$ are the length of the perpendiculars from origin on the tangent and normal drawn to the curve $x^{2/3}+y^{2/3}=6^{2/3}$, then find the value of $\dfrac{\left(4p_1^2+p_2^2\right)}{6}$.

\end{multicols}
\end{document}
